\section{Introduction}
\label{sec:introduction}

A finite reservoir changes the probabilities of subsystem energies.
Away from a phase transition, an energy window of width proportional
to $\sqrt N$ often contains nearly all canonical probability. At
first-order coexistence, two or more such windows can be separated by
an energy difference proportional to $N$. A reservoir that is adequate
within one window can then change the relative phase populations,
the fluctuations inside the phases, or both. The question addressed
here is how large a physical thermal reservoir must be to reproduce
the complete canonical law of an interacting subsystem, after its
total energy has been calibrated as favorably as possible.

The question arises wherever a small system exchanges energy with a
bath that is not much larger than itself. Finite-bath and Gaussian
ensemble Monte Carlo schemes deliberately use such baths to sample
first-order transitions
\cite{ChallaHetherington1988PRA,GriffinMattySwendsen2017}; nanoscale
calorimetry of clusters and confined fluids compares a finite sample
with a finite thermal environment; and a subsystem of a closed
classical or quantum system thermalizes only to the accuracy that the
remainder can supply \cite{RieraGogolinEisert2012}. In each case a
single number, the bath heat capacity, is expected to control the
error, and the practical question is how it must scale with the
subsystem size at coexistence.

We use total variation,
$\TV(P,Q)=\sup_A|P(A)-Q(A)|$.
It controls the error in every bounded measurement of the subsystem,
including its phase label and its fluctuation-scale observables.
Agreement of thermodynamic potentials, energy density, or a specified
finite set of averages does not in general establish full-law TV
convergence. Conversely, convergence of unbounded averages requires
additional integrability, as discussed in Section~\ref{sec:framework}.
Full-law TV is also a
different requirement from a vanishing unscaled Euclidean distance
between broad energy histograms. Because our reservoir likelihood
depends only on subsystem energy, energy-law TV equals full-state TV
exactly. Energy is therefore sufficient for computing the strong
error without replacing the microscopic system by an energy-only
model.

The reservoir has density of states
$\omega_B(U)\propto U^{c_N}\ind_{\{U>0\}}$.
The exponent $c_N>0$ is its surface-entropy heat capacity in units of
$k_B$; Section~\ref{sec:framework} states the finite-system entropy
convention and its distinction from the canonical heat capacity.
The subsystem Hamiltonian and target canonical inverse temperature
are fixed in each comparison. We optimize only the total energy of
the isolated composite. Weak interaction permits energy exchange but
contributes no interaction energy to this equilibrium marginal.

\subsection{The mechanism in two lines}

The subsystem sees the bath through the log likelihood
$\beta e+c_N\log(\mathcal E_N-e)$, a concave function of the
subsystem energy $e$ with curvature of order $1/c_N$. Choosing the
total energy $\mathcal E_N$ so that this function takes equal values at
two phase centers separated by $\Delta_N$ leaves endpoint slopes of
order $\Delta_N/c_N$. Across a phase of energy width $s_N$ such a slope
tilts the log weight by $\Delta_Ns_N/c_N$, so the full law survives
only if $c_N\gg\Delta_Ns_N$; for an extensive latent heat and
$\sqrt N$ fluctuations this is $c_N\gg N^{3/2}$. If instead three
energies on a common scale $b_N$ must all receive equal weight, no
choice of $\mathcal E_N$ helps: a concave function of curvature $1/c_N$
cannot be flat at three points spaced by $b_N$ unless $b_N^2/c_N\to0$.
The rest of the paper makes both statements sharp, in both directions,
and verifies their hypotheses for interacting spin models.

\subsection{Results}

The required capacity depends on the energy support that must be
preserved. If $(E_N-a_N)/b_N$ has a weak limit containing at least
three support points, the optimal condition is
$c_N\gg b_N^2$. This theorem requires no density, central limit
theorem, or moment bound. Here $x_N\gg y_N$ means $x_N/y_N\to\infty$.
For a single phase with a nondegenerate Gaussian limit on scale
$\sqrt N$, it gives $c_N\gg N$. For three distinct macroscopic energy
phases it gives $c_N\gg N^2$.

Two energy phases admit the secant calibration just described.
We prove the condition $c_N\gg\Delta_Ns_N$
necessary when at least one phase has a nondegenerate weak fluctuation
limit, and sufficient under explicit positive-phase exponential-moment
and exceptional-mass bounds. Only two distinct
\emph{energies} are involved: several symmetry-related ordered phases
at the same energy do not create a three-energy obstruction.

\begin{table}[t]
\centering
\caption{Optimal physical reservoir scales. Each convergence statement
concerns full-state TV and permits total-energy calibration.
The table summarizes asymptotic theorems under the indicated hypotheses;
it is not a universal finite-size error bound.}
\label{tab:threshold-summary}
\small
\renewcommand{\arraystretch}{1.12}
\begin{tabular}{@{}>{\raggedright\arraybackslash}p{.29\linewidth}>{\raggedright\arraybackslash}p{.21\linewidth}>{\raggedright\arraybackslash}p{.41\linewidth}@{}}
\toprule
Canonical energy structure & Required capacity & Hypotheses and result\\
\midrule
One ordinary fluctuation window & $c_N\gg N$ &
Nondegenerate weak Gaussian limit on scale $\sqrt N$;
Theorem~\ref{thm:three-support}\\[3pt]
Two extensive energy phases & $c_N\gg N^{3/2}$ &
Positive phases, one nondegenerate $\sqrt N$ limit, controlled
phase tails and exceptional mass;
Corollary~\ref{cor:two-phase-iff}\\[3pt]
At least three extensive energy phases & $c_N\gg N^2$ &
Weak macroscopic limit with at least three support points;
Theorem~\ref{thm:three-support}\\[3pt]
Two independent two-phase copies, joint target & $c_N\gg N^2$ &
Three support points for the total energy;
Corollary~\ref{cor:shared-joint-threshold}\\
\bottomrule
\end{tabular}
\end{table}

The two-phase theorem is realized by both a three-state
Curie--Weiss--Potts Hamiltonian and a short-range Potts Hamiltonian
with a fixed large number of colors, in every fixed spatial dimension
$d\ge2$. Quadratic kinetic coordinates are optional in both models.
For the short-range system, the main difficulty is controlling the
positive probability carried by interfacial configurations after
reservoir reweighting. We obtain phase-energy exponential moments
from positive restricted contour partition functions, then transfer
them from occupied bonds to the exchanged spin energy. A signed
finite-volume partition error alone would not establish this control.

For these microscopic examples, let $w_-,w_+$ denote the limiting
phase probabilities. The subsystem's canonical heat capacity is itself
of order $N^2$ at coexistence:
\begin{equation}
 \frac{C_{S,\mathrm{can}}}{k_BN^2}
 =\frac{\beta_N^2\Var_{P_N}(E_N)}{N^2}
 \longrightarrow\beta^2w_-w_+\ell^2,
 \qquad \ell=\lim_N\Delta_N/N.
 \label{eq:canonical-coexistence-capacity}
\end{equation}
Indeed, the bounded spin energy per particle converges to the two-atom
law, so its first two moments converge. Optional independent kinetic
energy with Gamma shape $aN$ adds only $aN/\beta_N^2$ to the variance;
$a=0$ denotes no kinetic sector. The proved phase
moments bound within-phase variances by order $N$; the leading
$N^2$ contribution instead comes from the separation of the phases.
Secant calibration compensates their relative bias, permitting full-law
accuracy with $N^{3/2}\ll c_N\ll N^2$, even though the bath heat
capacity is then smaller than this full canonical coexistence heat
capacity. This is the physical content of the $N^{3/2}$ scale: a bath
need not be able to absorb the latent heat at a fixed temperature in
order to reproduce the coexisting law exactly.

A shared bath reveals a second distinction. For two otherwise
independent copies, a capacity in the window $N\ll c_N\ll N^2$ can
suppress equal-phase pairs while retaining ordinary fluctuations in
opposite-phase pairs. At balanced canonical phase weights, both full
subsystem marginals converge to their canonical laws, whereas joint
TV tends to $1/2$ and the full microscopic mutual information tends
to $\log2$. The reference temperature needed to balance the phases
is specified and used consistently by the bath; for the spin models
it is a shift of order $1/N$ from the coexistence temperature. Joint
canonical accuracy requires the larger $N^2$ scale because the sum of
two two-phase energies has three macroscopic support points.

At the boundary $c_N\sim\gamma N^{3/2}$, stronger phase-tail
integrability and exceptional-mass suppression allow us to determine
the finite Gaussian phase tilts, their integrated probabilities, and
the full TV error. These hypotheses hold for every fixed $\gamma>0$
in the mean-field model and the short-range model in dimensions above
two; in two dimensions our proof gives them for sufficiently large
$\gamma$. We optimize this error over every admissible composite
energy, including choices that discard a phase. This matters
operationally: minimizing one scalar distance can favor losing a
phase rather than retaining its probability with distorted
fluctuations. Exact multivariate Gaussian models in
Appendix~\ref{sec:gaussian-geometry} show a related geometric
constraint: the phase points that can survive optimal quadratic
reweighting are governed by empty-sphere contact sets.

\subsection{Calibration sensitivity}

The secant calibration must be accurate. In the regime
$c_N\gg\Delta_N$, shifting $\mathcal E_N$ by $\delta$ changes the
difference of the two center log weights by approximately
$\beta_N^2\Delta_N\delta/c_N$, so the phase odds change by a factor
of order one when $|\delta|\asymp c_N/(\beta_N^2\Delta_N)$.
Preserving the limiting phase weights of the calibrated law instead
requires $\delta=o(c_N/(\beta_N^2\Delta_N))$.
At either center the bath energy is
$U_{B,i,N}=\mathcal E_N-e_{i,N}\sim c_N/\beta_N$.
Thus, for the microscopic models with $\Delta_N\asymp N$, relative
bath-energy errors of order $1/N$ change the phase odds by order one;
vanishing changes require relative errors $o(1/N)$. The same holds
for relative errors in the bath temperature at either center.
The canonical ensemble has exactly the same sensitivity at
coexistence: a change of order $1/N$ in $\beta$ multiplies the phase
odds by a factor of order one. The requirement is therefore a property
of the target law, not of the finite bath. Section~\ref{sec:boundary-smooth}
quantifies the effect of calibration errors of order $\sqrt N$ at the
$N^{3/2}$ boundary, where they change the limiting phase weights but
not the fluctuation tilts.

\subsection{Relation to earlier work}

Finite-reservoir ensembles and their fluctuation corrections are
established. Challa and Hetherington introduced Gaussian ensembles
as a controllable interpolation between familiar ensembles and
demonstrated changes to first-order energy histograms
\cite{ChallaHetherington1988PRL,ChallaHetherington1988PRA}.
The two-Gaussian coexistence description of Challa, Landau, and
Binder \cite{ChallaLandauBinder1986} already separates an extensive
latent-energy gap from ordinary within-phase widths, and their later
review covers finite-size and finite-bath approaches
\cite{ChallaLandauBinder1990}.
Power-law subsystem laws from finite constant-capacity reservoirs
are standard \cite{Campisi2007}, and the surface-versus-volume entropy
convention for the bath is the one compared in
\cite{HilbertHanggiDunkel2014,DunkelHilbert2014}.

Strong-distance bath bounds also exist.
Riera, Gogolin, and Eisert \cite{RieraGogolinEisert2012} control a
reduced microcanonical-shell state in trace distance by the
nonlinear remainder of the bath entropy; in their bath example a
sufficient size is proportional to the square of the subsystem
energy range, which in a commuting setting is a TV bound of the
$N^2$ type. Sharp growing-block conditional limit theorems
\cite{DiaconisFreedman1988} treat a regular exponential-family
conditioning problem rather than the full phase mixture of an
interacting subsystem coupled to a separate reservoir.
Finite-bath Potts comparisons by Griffin, Matty, and Swendsen
\cite{GriffinMattySwendsen2017} are the closest numerical precedent;
they measure agreement by an unweighted Euclidean histogram norm and
optimize a comparison temperature, and
Appendix~\ref{sec:capillarity-diagnostics} gives an exact example
in which that norm vanishes while full-state TV does not.
Generalized Gaussian ensembles recover nonconcave entropy and
macrostate equivalence at fixed penalty
\cite{CosteniucEllisTouchette2006,CosteniucEllisTouchetteTurkington2005},
and squeezed ensembles select first-order states with finite-size
conversion formulas \cite{YonetaShimizu2019}; both are
thermodynamic-limit questions at fixed penalty rather than
complete-law asymptotics at shrinking curvature.
For shared baths, ensemble dependence of fluctuations and
correlations is classical \cite{LebowitzPercusVerlet1967,Mishin2015};
opposite-phase assignments in two weakly coupled identical systems
were obtained by Ram\'irez-Hern\'andez, Larralde, and Leyvraz
\cite{RamirezHernandezLarraldeLeyvraz2008}; and ensemble-dependent
shared information, including order-one contributions from
degenerate phases, was studied by Cohen, Rittenberg, and Sadhu
\cite{CohenRittenbergSadhu2015}.

Against this background, the contributions of the present paper are
the following. First, sharp necessity: for the exact power-law bath
and every total-energy calibration, a weak energy limit with three
support points forces $c_N\gg b_N^2$, and two phases with a
nondegenerate fluctuation limit force $c_N\gg\Delta_Ns_N$. Second,
the optimized two-phase scale $N^{3/2}$ is attained, below both the
energy-range-squared scale and the coexistence heat capacity, under
positive phase-tail hypotheses that we state explicitly. Third, these
hypotheses are verified for mean-field and short-range Potts
Hamiltonians, the latter through positive contour measures and a
bond-to-energy transfer. Fourth, the boundary law at
$c_N\asymp N^{3/2}$ is computed and optimized over every composite
energy, including phase abandonment. Fifth, under a shared bath the
full marginal, joint, and information limits are proved in a specified
capacity window.

\subsection{Organization}

Section~\ref{sec:framework} derives the exact reservoir likelihood
and calibrations. Section~\ref{sec:thresholds} proves the weak-support
and two-phase criteria. Section~\ref{sec:microscopic} gives
microscopic realizations and exact finite-system formulas;
Appendix~\ref{app:short-range-proof} contains the contour proof.
Section~\ref{sec:shared-baths} treats joint statistics under a
shared bath, and Section~\ref{sec:boundary-smooth} gives smooth-bath
extensions, the optimized physical boundary, and the exact interior
gain of the secant bath.
Section~\ref{sec:numerics} presents reproducible deterministic
calculations. The exact Gaussian geometry, the capillarity and
accuracy diagnostics, and the shared-bath extensions to fixed phase
counts and linear capacity are collected in
Appendices~\ref{sec:gaussian-geometry}, \ref{sec:capillarity-diagnostics},
and~\ref{app:shared-extensions}.
